% Original minimal reproductions derived from the corpus review.
% Each frame is isolated by scripts/probe-beamer-frame.mjs. Keep all settings
% needed by the oracle in the preamble or within the selected frame.
\documentclass[aspectratio=169]{beamer}
\usepackage{amsmath}
\usepackage{booktabs}
\usepackage{graphicx}
\usetheme{Madrid}
\usecolortheme{seahorse}
\title{Corpus priorities}
\author{}
\date{}

\begin{document}

% 1: positive control, identical flow shape to hello_world_beamer.tex.
\begin{frame}{Plain control}
  \begin{columns}[T,totalwidth=\textwidth]
    \begin{column}{\textwidth}Baseline body text.\end{column}
  \end{columns}
\end{frame}

% 2: supported siblings and nested lists within a block/column.
\begin{frame}{Nested flow}
  \begin{columns}[T,totalwidth=\textwidth]
    \begin{column}{.47\textwidth}
      Before the block.
      \begin{block}{Result}
        \begin{itemize}
          \item Alpha
          \begin{itemize}\item Beta\end{itemize}
          \item Gamma
        \end{itemize}
      \end{block}
      After the block.
    \end{column}
    \begin{column}{.47\textwidth}Other column.\end{column}
  \end{columns}
\end{frame}

% 3: remove the unsupported listing from r2-43 to isolate the block's math.
\begin{frame}{Block inline math}
  \begin{exampleblock}{Result}
    Alpha $a\cdot b=\sum_{i=1}^{n}a_i b_i$.
  \end{exampleblock}
\end{frame}

% 4: r2-34's block body, with neither custom macros nor theme overrides.
\begin{frame}{Block aligned matrix}
  Before the block.
  \begin{block}{Encoding}
    \begin{align*}
      \mathrm{Alpha}&=\begin{pmatrix}1&0&0\end{pmatrix}\\
      \mathrm{Beta}&=\begin{pmatrix}0&1&0\end{pmatrix}
    \end{align*}
  \end{block}
  After the block.
\end{frame}

% 5: display math must not swallow the surrounding list.
\begin{frame}{List with cases}
  \begin{itemize}
    \item Alpha
    \item The function is
      \[f(x)=\begin{cases}0&x<0\\x&x\geq0\end{cases}\]
    \item Beta
  \end{itemize}
\end{frame}

% 6: r2-17's leading footnote plus two ordinary lists.
\begin{frame}{Footnote and lists}
  \footnote{Contact: \url{https://example.org}.}
  Alpha
  \begin{itemize}\item Beta\item Gamma\end{itemize}
  Delta
  \begin{itemize}\item Epsilon\end{itemize}
\end{frame}

% 7/8: same asset through a figure wrapper and directly. Keep its natural
% TeX size (64bp = 64.24pt) so the Lua node trace needs no PDF image matrix.
\begin{frame}{Figure caption}
  Before the figure.
  \begin{figure}
    \centering
    \includegraphics[width=64.24pt]{test-image.png}
    \caption{Alpha image}
  \end{figure}
  After the figure.
\end{frame}

\begin{frame}{Direct image control}
  \includegraphics[width=64.24pt]{test-image.png}

  Alpha image.
\end{frame}

% 9: simplest table, no package-specific column types or wrapper.
\begin{frame}{Basic table}
  \begin{tabular}{lr}
    Alpha & 1\\
    Beta & 2\\
    Gamma & 3
  \end{tabular}
\end{frame}

% 10: booktabs rules and a table wrapper in a column.
\begin{frame}{Table in columns}
  \begin{columns}[T,totalwidth=\textwidth]
    \begin{column}{.47\textwidth}
      \begin{table}
        \begin{tabular}{lr}
          \toprule
          Alpha & Value\\
          \midrule
          Beta & 2\\
          Gamma & 3\\
          \bottomrule
        \end{tabular}
      \end{table}
    \end{column}
    \begin{column}{.47\textwidth}Other column.\end{column}
  \end{columns}
\end{frame}

% 11: isolate rotation from the table in r2-21.
\begin{frame}{Rotated text}
  Before.
  \rotatebox{90}{Alpha}
  After.
\end{frame}

% 12: scaling and minipage composition, without a table or author macro.
\begin{frame}{Scaled minipage}
  \scalebox{.8}{\begin{minipage}{160pt}
    Alpha paragraph.

    Beta paragraph.
  \end{minipage}}
\end{frame}

% 13: declaration-form onslide scope extends until the next declaration.
\begin{frame}{Declaration overlays}
  \onslide<1->Alpha\par
  \onslide<2->Beta\par
  \onslide<3->Gamma\par
\end{frame}

% 14: built-in overprint and handout selectors from r2-29.
\begin{frame}{Overprint states}
  \begin{overprint}
    \onslide<1|handout:0>Alpha.
    \onslide<2|handout:0>Beta.\par Second line.
    \onslide<3|handout:1>Gamma.
  \end{overprint}
\end{frame}

% 15: positive control for explicit list overlays, including the hidden state.
\begin{frame}{Item overlay control}
  \begin{itemize}
    \item<1-> Alpha
    \item<2-> Beta
  \end{itemize}
\end{frame}

% 16: three continuation pages are not three overlay states.
\begin{frame}[allowframebreaks]{Forced continuations}
  Alpha.
  \newpage
  Beta.\par
  \framebreak
  Gamma.
\end{frame}

% 17: automatic continuation without a forced page break.
\begin{frame}[allowframebreaks=1]{Automatic continuations}
  \begin{itemize}
    \item \rule{0pt}{30pt}Alpha
    \item \rule{0pt}{30pt}Beta
    \item \rule{0pt}{30pt}Gamma
    \item \rule{0pt}{30pt}Delta
    \item \rule{0pt}{30pt}Epsilon
    \item \rule{0pt}{30pt}Zeta
    \item \rule{0pt}{30pt}Eta
    \item \rule{0pt}{30pt}Theta
  \end{itemize}
\end{frame}

% 18: equation+split with a local reference, all inside the isolated frame.
\begin{frame}{Split equation and reference}
  \begin{equation}\label{eq:priority}
    \begin{split}
      a+b&=c\\
      d+e&=f
    \end{split}
  \end{equation}
  Alpha equation \ref{eq:priority}.
  \begin{equation}\begin{split}x+y&=z\\u+v&=w\end{split}\end{equation}
\end{frame}

% 19: cases in a column, r2-03 with the unrelated plot removed.
\begin{frame}{Cases in columns}
  \begin{columns}[T,totalwidth=\textwidth]
    \begin{column}{.47\textwidth}
      \[f(x)=\begin{cases}0&\text{if }x<0\\x&\text{if }x\geq0\end{cases}\]
      Alpha.
    \end{column}
    \begin{column}{.47\textwidth}Other column.\end{column}
  \end{columns}
\end{frame}

% 20: ordinary scoped sizing and list spacing, not opaque code syntax.
\begin{frame}{Scoped size and list spacing}
  \begin{small}
    Alpha paragraph.
    \begin{itemize}
      \setlength{\itemsep}{1em}
      \item Beta
      \item Gamma
    \end{itemize}
  \end{small}
  Delta paragraph.
\end{frame}

\end{document}
